% ============================================================
\section{Definição de Qualidade de Serviço}
% ============================================================

A Qualidade de Serviço (QoS) é o uso de mecanismos e tecnologias que atuam em
uma rede com o objetivo de controlar o tráfego e assegurar a boa
\textit{performance} de aplicações críticas mesmo diante de capacidade de rede
limitada \cite{fortinet2024}. Isso permite que organizações ajustem o tráfego de
rede priorizando o desempenho das aplicações mais importantes.

\subsection{Funcionamento Básico}

Quando computadores enviam e recebem informações em uma rede, essas informações
são organizadas em \textbf{pacotes}. O QoS funciona \textbf{marcando pacotes}
para identificar os tipos de serviço e, em seguida, \textbf{configurando
roteadores} para criar filas virtuais separadas para cada aplicação, baseadas em
prioridade \cite{fortinet2024}. Como resultado, largura de banda é reservada para
aplicações críticas ou de maior prioridade na fila.

Uma ferramenta de QoS observa o cabeçalho (\textit{header}) do pacote para
decidir quais pacotes priorizar. O cabeçalho contém informações sobre o tipo de
dado transportado, o destino na rede e a finalidade do pacote
\cite{techtarget2024}.

\subsection{Serviços que Requerem QoS}

O QoS é comumente aplicado em redes que possuem tráfego de sistemas que requerem
alta quantidade de recursos. Entre os principais serviços que se beneficiam do
QoS, destacam-se \cite{fortinet2024}:

\begin{itemize}
  \item \textit{Internet Protocol Television} (IPTV)
  \item \textit{Online gaming}
  \item \textit{Streaming} de mídia
  \item Videoconferência
  \item \textit{Video on Demand} (VOD)
  \item \textit{Voice over IP} (VoIP)
\end{itemize}

% ============================================================
\section{Elementos Avaliados}
% ============================================================

Para garantir a qualidade de serviço, é necessário monitorar e controlar um
conjunto de parâmetros de rede. A seguir são descritos os principais elementos
avaliados pelo QoS \cite{geeksforgeeks2024}.

\subsection{Largura de Banda}

A \textbf{largura de banda} (\textit{bandwidth}) é a capacidade máxima que um
enlace de rede possui para transmitir a maior quantidade de dados de um ponto a
outro no menor tempo possível \cite{fortinet2024}. O QoS instrui o roteador sobre
como utilizar essa largura de banda — por exemplo, alocando determinada
quantidade para diferentes filas por tipo de tráfego.

\subsection{Atraso}

O \textbf{atraso} (\textit{delay} ou latência) é o tempo que um pacote leva da
sua origem até o destino \cite{fortinet2024}. Pode ser afetado pelo
\textit{queuing delay} (atraso na fila), que ocorre em períodos de
congestionamento. O QoS reduz o atraso criando filas de prioridade para certos
tipos de tráfego. A latência idealmente deve ser a menor possível; em serviços de
VoIP, por exemplo, muita latência pode causar sobreposição dos áudios durante uma
conversa.

\subsection{Variação de Atraso (Jitter)}

O \textbf{jitter} (variação de atraso) é a irregularidade na velocidade de
chegada dos pacotes na rede, resultante do congestionamento \cite{fortinet2024}.
Pacotes podem chegar atrasados ou fora de ordem, gerando quebras em áudios e
vídeos transmitidos.

\subsubsection{Uso de Buffer para Redução de Jitter}

Uma técnica amplamente utilizada para mitigar o jitter é o emprego de
\textbf{buffers de jitter} (também chamados de \textit{playout buffers}). Esses
buffers são áreas de armazenamento temporário que coletam os pacotes recebidos,
armazenam-nos por um breve período e os encaminham ao processador em intervalos
uniformes, compensando a variação de atraso antes da reprodução.

Existem dois tipos principais de buffer de jitter:

\begin{itemize}
  \item \textbf{Buffer estático}: possui tamanho fixo, configurado manualmente
  pelo administrador. Simples de implementar, mas não se adapta a variações nas
  condições da rede.
  \item \textbf{Buffer adaptativo}: ajusta dinamicamente seu tamanho com base nas
  condições atuais da rede. Mais eficiente em ambientes com jitter variável.
\end{itemize}

O uso de buffer introduz uma latência constante e previsível, que é preferível à
variação imprevisível causada pelo jitter. Em aplicações VoIP, por exemplo, o
buffer tipicamente retém entre 50 ms e 150 ms de áudio para absorver as variações
de chegada dos pacotes sem causar interrupções perceptíveis na chamada.

\subsection{Perda de Pacotes}

A \textbf{perda de pacotes} (\textit{packet loss}) é a quantidade de dados
perdidos por conta de pacotes descartados, normalmente causada por congestionamento
na rede \cite{fortinet2024}. Quando o tráfego excede a capacidade dos roteadores e
switches, estes começam a descartar pacotes para conseguir lidar com a carga.

O QoS permite que as organizações definam \textbf{quais pacotes descartar} quando
isso ocorre, priorizando o tráfego crítico. Os pacotes geralmente são descartados
quando a fila de espera ultrapassa seu limite máximo. Em comunicações em tempo
real — seja vídeo ou áudio — a perda de pacotes pode causar jitter e travamentos
na fala \cite{fortinet2024}.

\subsection{Outras Métricas}

Além dos parâmetros principais, o QoS também avalia \cite{fortinet2024,techtarget2024}:

\begin{itemize}
  \item \textbf{Avaliação Média de Opinião (MOS)}: métrica para medir a qualidade
  do áudio usando escala de 1 a 5, onde 5 representa a melhor qualidade.
  \item \textbf{Capacidade de Processamento (\textit{Throughput})}: medida do
  sucesso na entrega de dados na rede.
  \item \textbf{\textit{Error Rate}}: mede a frequência de arquivos corrompidos e
  perdas de pacote durante a transmissão.
\end{itemize}

% ============================================================
\section{Técnicas de QoS}
% ============================================================

Diversas técnicas são empregadas para implementar QoS em redes modernas
\cite{fortinet2024,techtarget2024}.

\subsection{Marcação de Tráfego (\textit{Traffic Marking})}

O tráfego prioritário precisa ser identificado antes de ser tratado. A marcação
é realizada por meio de:

\begin{itemize}
  \item \textbf{CoS (\textit{Class of Service})}: marca o fluxo no cabeçalho do
  frame de camada 2.
  \item \textbf{DSCP (\textit{Differentiated Services Code Point})}: marca o fluxo
  no cabeçalho do pacote de camada 3.
\end{itemize}

\subsection{Filas (\textit{Queuing})}

O processo de \textit{queuing} cria políticas que fornecem tratamento preferencial
a certos fluxos. As filas são buffers de alta \textit{performance} em roteadores e
switches \cite{fortinet2024}. Pacotes com prioridade mais alta são movidos para uma
fila dedicada que transmite dados mais rapidamente, reduzindo as chances de
descarte. O \textit{priority queuing} garante disponibilidade e mínima latência para
as aplicações mais importantes.

\subsection{RSVP (\textit{Resource Reservation Protocol})}

O RSVP é um protocolo de sinalização que opera na camada de rede e reserva recursos ao longo de
uma rede para garantir níveis específicos de QoS para fluxos de dados de
aplicações \cite{fortinet2024}. Permite que aplicações solicitem uma reserva de
largura de banda antes de iniciar a transmissão.

\subsection{\textit{Traffic Shaping}}

O \textit{traffic shaping} é uma técnica de limitação de taxa que otimiza a
\textit{performance} e aumenta a largura de banda utilizável \cite{fortinet2024}.
Regula o fluxo de pacotes na saída, suavizando picos de tráfego e evitando
congestionamento.

\subsection{Abordagens de QoS}

Dois modelos principais são utilizados para gerenciar QoS em redes IP
\cite{ieee_ipv6qos}:

\begin{itemize}
  \item \textbf{IntServ (\textit{Integrated Services})}: fornece garantias de QoS
  por fluxo individual para sessões de aplicações, utilizando o protocolo RSVP
  para sinalizar e configurar o caminho. Sua principal desvantagem é a falta de
  escalabilidade, pois a carga de armazenamento e processamento nos roteadores
  aumenta proporcionalmente ao número de fluxos.
  \item \textbf{DiffServ (\textit{Differentiated Services})}: o tráfego é
  condicionado por um \textit{bandwidth broker} e classificado nas bordas da rede
  por meio do campo DSCP. Embora mais escalável, não consegue diferenciar fluxos
  individuais dentro de uma mesma classe.
\end{itemize}

Para implementar esses modelos, os componentes da rede são equipados com as
seguintes funcionalidades \cite{ieee_ipv6qos}:

\begin{itemize}
  \item \textbf{Controle de Admissão (\textit{Admission Control})}: determina se
  ainda há recursos disponíveis para um serviço solicitado e se eles podem ser
  fornecidos em uma interface específica.
  \item \textbf{Policiamento e Modelagem (\textit{Policing and Shaping})}: conjunto
  de ações realizadas quando o tráfego real excede os valores negociados. Incluem
  descarte de pacotes, rebaixamento para classe inferior ou marcação como não
  conforme.
  \item \textbf{Classificador de Pacotes (\textit{Packet Classifier})}: identifica
  pacotes de um fluxo particular e os associa a uma classe de QoS. Implementado
  nos pontos de borda (\textit{edge points}) da rede.
  \item \textbf{Agendador de Pacotes (\textit{Packet Scheduler})}: garante que os
  fluxos identificados recebam a Qualidade de Serviço solicitada. Também
  implementado nos pontos de borda.
\end{itemize}

% ============================================================
\section{QoS em Redes IPv6}
% ============================================================

O IPv6 introduz melhorias nativas para o QoS por meio de dois campos específicos
no cabeçalho do pacote \cite{ieee_ipv6qos}:

\begin{itemize}
  \item \textbf{Traffic Class (TC)}: utilizado para diferenciar prioridades de
  tráfego, de forma análoga ao campo DSCP do IPv4.
  \item \textbf{Flow Label}: identificador de 20 \textit{bits} que distingue
  fluxos originados da mesma fonte de forma exclusiva.
\end{itemize}

A combinação do endereço IP de origem com o \textit{Flow Label} resulta em um
tempo menor de busca nas tabelas dos roteadores, minimizando o atraso de
processamento.

\subsection{Flow Label}

Atribuir um valor de \textit{Flow Label} aos pacotes permite que a rede os
classifique utilizando apenas a semântica de IP, agilizando o processamento
\cite{ieee_ipv6qos}. Isso evita o problema da \textbf{Violação de Camada}
(\textit{Layer Violation}) — a prática de extrair dados das camadas de aplicação
superiores para processar pacotes. Evitar essa violação é essencial quando os
pacotes trafegam criptografados ou com os números de porta ocultos.

\subsection{Controle de Admissão em IPv6}

O controle de admissão é o algoritmo que decide se um novo fluxo pode ser
atendido pela rede sem prejudicar as reservas já existentes. Avalia a carga atual,
os recursos solicitados e as regras da política do serviço \cite{ieee_ipv6qos}.

\subsubsection{Bandwidth Broker}

O \textit{Bandwidth Broker} é um agente que aloca serviços para usuários e
configura os roteadores. Ele consulta um banco de dados de políticas, verifica se
há largura de banda suficiente, aprova a requisição do fluxo e configura o
roteador para aceitar o pacote \cite{ieee_ipv6qos}.

\subsubsection{QoSbox}

O QoSbox é um roteador configurável que adapta ativamente suas decisões de
policiamento e agendamento de acordo com a chegada instantânea do tráfego. Ao
contrário do \textit{Bandwidth Broker}, não depende de agentes externos: utiliza
\textit{buffers} por classe e algoritmos locais para gerenciar atrasos e descartar
pacotes quando necessário, mantendo o nível de serviço das demais classes
\cite{ieee_ipv6qos}.

% ============================================================
\section{QoS e Multimídia}
% ============================================================

A qualidade de serviço é essencial para serviços multimídia porque o tráfego de
áudio e vídeo é muito sensível ao atraso, variação de atraso (jitter), perda de
pacotes e à largura de banda disponível \cite{geeksforgeeks2024}. Qualquer
degradação nesses parâmetros resulta em quebras, travamentos e áudio
ininteligível.

O QoS contribui para aplicações multimídia das seguintes formas
\cite{fortinet2024,geeksforgeeks2024}:

\begin{itemize}
  \item Prioriza serviços como VoIP, chamadas de vídeo e \textit{live stream}
  para evitar interrupções.
  \item Controla a latência e o jitter para manter áudio e vídeo sincronizados.
  \item Reduz a perda de pacotes priorizando pacotes de mídia durante períodos de
  congestionamento.
  \item Aloca mais largura de banda para fluxos de dados multimídia, permitindo
  maior \textit{bitrate} sem prejudicar outros tráfegos críticos.
  \item Otimiza o uso de recursos ao agendar e enfileirar pacotes, permitindo que
  dados multimídia e outros dados compartilhem a rede de forma eficiente.
\end{itemize}

% ============================================================
\section{Forward Error Correction (FEC)}
% ============================================================

O \textit{Forward Error Correction} (FEC) é uma técnica que fornece ao receptor
a capacidade de corrigir erros e reproduzir pacotes de forma imediata
\cite{geeksforgeeksfec2024}. O conceito central é a inserção prévia de redundância
matemática nos dados enviados, eliminando a necessidade de um canal reverso para
solicitar a retransmissão de informações perdidas ou corrompidas.

Três técnicas principais de FEC são utilizadas \cite{geeksforgeeksfec2024}:

\begin{itemize}
  \item \textbf{Distância de Hamming}: método matemático onde a distância mínima
  para corrigir $t$ erros é definida por $d_{min} = 2t+1$. Raramente usada em
  redes por exigir uma quantidade inviável de \textit{bits} redundantes.
  \item \textbf{XOR}: recria pacotes perdidos por meio da propriedade
  ou-exclusivo. Um pacote é dividido em $N$ fragmentos; a operação XOR é calculada
  entre eles e a rede envia $N+1$ fragmentos. Para $N=4$, envia-se 25\% de dados
  adicionais, permitindo corrigir a perda de um dos quatro fragmentos.
  \item \textbf{Entrelaçamento de Fragmentos (\textit{Chunk Interleaving})}: cada
  pacote é dividido em pequenos blocos criados horizontalmente, mas enviados
  verticalmente. Cada pacote transmitido carrega um pequeno pedaço de vários
  pacotes originais diferentes.
\end{itemize}

\subsection{Exemplo de Uso}

O FEC é amplamente utilizado em comunicações multimídia de tempo real
\cite{geeksforgeeksfec2024}. Com a técnica de \textit{Chunk Interleaving}, se um
pacote for completamente perdido durante a transmissão, o receptor perderá apenas
um fragmento muito pequeno de cada pacote original. Como pequenos fragmentos
ausentes são normalmente aceitáveis em fluxos de áudio ou vídeo, a reprodução
continua sem interrupções perceptíveis.

\subsection{Relação com QoS}

A retransmissão de pacotes perdidos é inviável em aplicações de tempo real, pois
gera um atraso inaceitável na reprodução \cite{geeksforgeeksfec2024}. A principal
relação do FEC com o QoS está na capacidade de proteger os parâmetros de
\textbf{Atraso} e \textbf{Variação de Atraso (Jitter)}: ao reproduzir pacotes
perdidos diretamente no destino de forma imediata, o FEC preserva a fluidez da
aplicação e mantém os requisitos de tráfego, evitando a degradação da Qualidade
de Serviço.

% ============================================================
\section{Qualidade de Experiência (QoE)}
% ============================================================

A \textbf{Qualidade de Experiência} (QoE, do inglês \textit{Quality of
Experience}) é uma medida do nível geral de satisfação e experiência de um
usuário com um produto ou serviço \cite{techtargetqoe2024}. A União Internacional
de Telecomunicações define QoE como a aceitabilidade geral de uma aplicação ou
serviço, conforme percebida subjetivamente pelo usuário final.

Embora relacionados, QoS e QoE não são o mesmo conceito. O QoS refere-se às
métricas técnicas e objetivas da rede — atraso, jitter, perda de pacotes —
enquanto o QoE adota a perspectiva do usuário final, buscando responder à
pergunta: \textit{``Este serviço entregou uma experiência boa o suficiente?''}
\cite{techtargetqoe2024}.

O método mais comum para avaliar o QoE é pesquisar ou amostrar um grande número
de usuários, coletando percepções subjetivas sobre a qualidade do serviço
recebido.

% ============================================================
\section{Vantagens e Melhores Práticas}
% ============================================================

\subsection{Vantagens do QoS}

A adoção de QoS traz diversos benefícios às organizações \cite{fortinet2024}:

\begin{enumerate}
  \item \textbf{Priorização de aplicações}: garante que as aplicações mais
  críticas sempre terão os recursos necessários para atingir alta
  \textit{performance}.
  \item \textbf{Melhor gerenciamento de recursos}: permite que administradores
  gerenciem melhor os recursos de rede, reduzindo custos e a necessidade de
  expansões de enlace.
  \item \textbf{Melhora da experiência do usuário}: o objetivo final é garantir a
  alta \textit{performance} das aplicações críticas, culminando em maior
  produtividade.
  \item \textbf{Gerenciamento de tráfego ponto a ponto}: permite entregar pacotes
  em ordem de um ponto a outro na internet, sem sofrer perda de pacotes.
  \item \textbf{Prevenção de perda de pacotes}: o QoS evita a perda priorizando a
  largura de banda para aplicações de alta \textit{performance}.
  \item \textbf{Redução de latência}: o QoS reduz a latência ao priorizar as
  aplicações críticas na fila de processamento.
\end{enumerate}

\subsection{Melhores Práticas}

Para uma configuração eficaz de QoS, recomenda-se \cite{fortinet2024}:

\begin{enumerate}
  \item Não definir limites de largura de banda máxima muito baixos na interface
  de origem, para evitar descarte excessivo de pacotes.
  \item Considerar a razão de distribuição de pacotes entre as filas disponíveis e
  quais filas são usadas por quais serviços, pois isso afeta latência e
  distribuição de fila.
  \item Aplicar garantias de largura de banda apenas em serviços específicos,
  evitando que todo o tráfego use a mesma fila em situações de alto volume.
  \item Configurar a priorização por \textbf{um único tipo}: ou por prioridade
  baseada em serviço, ou por prioridade de política de segurança — não ambos.
  \item Minimizar a complexidade da configuração de QoS para garantir alta
  \textit{performance}.
  \item Para resultados de teste precisos, usar o UDP e não superalocar o
  \textit{throughput} de largura de banda.
\end{enumerate}
